% This LaTeX file was created by Metamath on 22-Jul-2024 3:28 PM.
\begin{restatable}{axiom}{axbltaut}~\blTitle{ax-bl.taut}~ Tautologies hold in interpretation
$\mathbf{I}$.
   ax-bl.taut says that if $\mw{\varphi}$ is tautology, it also holds in
interpretation $\mathbf{I}$.
\begin{align}
\vdash ( \mw{\varphi} \rightarrow \mathbf{I} \lsem \mw{\varphi} \rsem
    )\label{eq:ax-bl.taut}\tag{ax-bl.taut}
\end{align}
\end{restatable}

\begin{restatable}{lemma}{blty}~\blTitle{bl.ty}~ Set membership in interpretation $\mathbf{I}$.
Types of variables are constant under interpretation $\mathbf{I}$.  (Unproven
in SFS.mm as of this writing.)
\begin{align}
\vdash ( \mc{A} \in \mc{B} \rightarrow \mathbf{I} \lsem \mc{A} \in \mc{B}
    \rsem )\label{eq:bl.ty}\tag{bl.ty}
\end{align}
\end{restatable}

\begin{restatable}{metadefinition}{dfblim}~\blTitle{df-bl.im}~ Interpretation of implication.  Only
the sound (modal K-axiom) direction is axiomatized: the converse is false
in general (Lean/Mathlib formalization, Lean4SFS/SFS.lean's dl\_im\_sound,
found a counterexample).
\begin{align}
\vdash ( \mathbf{I} \lsem ( \mw{\varphi} \rightarrow \mw{\psi} ) \rsem
    \rightarrow ( \mathbf{I} \lsem \mw{\varphi} \rsem \rightarrow \mathbf{I}
    \lsem \mw{\psi} \rsem ) )\label{eq:df-bl.im}\tag{df-bl.im}
\end{align}
\end{restatable}

\begin{restatable}{lemma}{blmono}~\blTitle{bl.mono}~ Monotonicity of $\mathbf{I}$: any provable
implication lifts through the interpretation.  This is the general tool for
weakening a plain-logic fact $( \mw{\varphi} \rightarrow \mw{\psi} )$ into a
$\mathbf{I}$-level implication, used repeatedly below wherever a
formerly-axiomatized converse direction (df-bl.or, df-bl.not, bl.3or,
bl.ex, ...) was removed as unsound.
\begin{align}
\vdash ( ( \mw{\varphi} \rightarrow \mw{\psi} ) \rightarrow ( \mathbf{I} \lsem
    \mw{\varphi} \rsem \rightarrow \mathbf{I} \lsem \mw{\psi} \rsem )
    )\label{eq:bl.mono}\tag{bl.mono}
\end{align}
\end{restatable}

\begin{proof}
\begin{align}
  1 &&  & \vdash ( ( \mw{\varphi} \rightarrow \mw{\psi} ) \rightarrow
      \mathbf{I} \lsem ( \mw{\varphi} \rightarrow \mw{\psi} ) \rsem
      )&(\mbox{\ref{eq:ax-bl.taut}})\notag
  \\ 2 &&  & \vdash ( \mathbf{I} \lsem ( \mw{\varphi} \rightarrow
      \mw{\psi} ) \rsem \rightarrow ( \mathbf{I} \lsem \mw{\varphi}
      \rsem \rightarrow \mathbf{I} \lsem \mw{\psi} \rsem )
      )&(\mbox{\ref{eq:df-bl.im}})\notag
  \\ 3 &&  & \vdash ( ( \mw{\varphi} \rightarrow \mw{\psi} )
      \rightarrow ( \mathbf{I} \lsem \mw{\varphi} \rsem \rightarrow
      \mathbf{I} \lsem \mw{\psi} \rsem )
      )&(\mbox{\ref{eq:syl}},1,2)\notag
\end{align}
\end{proof}

\begin{restatable}{metadefinition}{dfblbi}~\blTitle{df-bl.bi}~ Interpretation of biconditional.  Only
the sound direction is axiomatized: the converse is false in general
(Lean4SFS/SFS.lean's dl\_bi\_sound).
\begin{align}
\vdash ( \mathbf{I} \lsem ( \mw{\varphi} \leftrightarrow \mw{\psi} ) \rsem
    \rightarrow ( \mathbf{I} \lsem \mw{\varphi} \rsem \leftrightarrow
    \mathbf{I} \lsem \mw{\psi} \rsem ) )\label{eq:df-bl.bi}\tag{df-bl.bi}
\end{align}
\end{restatable}

\begin{restatable}{metadefinition}{dfblan}~\blTitle{df-bl.an}~ Interpretation of conjunction.
\begin{align}
\vdash ( \mathbf{I} \lsem ( \mw{\varphi} \wedge \mw{\psi} ) \rsem
    \leftrightarrow ( \mathbf{I} \lsem \mw{\varphi} \rsem \wedge \mathbf{I}
    \lsem \mw{\psi} \rsem ) )\label{eq:df-bl.an}\tag{df-bl.an}
\end{align}
\end{restatable}

\begin{restatable}{lemma}{bl3an}~\blTitle{bl.3an}~ Interpretation of conjunction, three wff.
\begin{align}
\vdash ( \mathbf{I} \lsem ( \mw{\varphi} \wedge \mw{\psi} \wedge \mw{\chi} )
    \rsem \leftrightarrow ( \mathbf{I} \lsem \mw{\varphi} \rsem \wedge
    \mathbf{I} \lsem \mw{\psi} \rsem \wedge \mathbf{I} \lsem \mw{\chi} \rsem )
    )\label{eq:bl.3an}\tag{bl.3an}
\end{align}
\end{restatable}

\begin{proof}
\begin{align}
  1 &&  & \vdash ( ( \mw{\varphi} \wedge \mw{\psi} \wedge \mw{\chi} )
      \leftrightarrow ( ( \mw{\varphi} \wedge \mw{\psi} )
      \wedge \mw{\chi} ) )&(\mbox{\ref{eq:df-3an}})\notag
  \\ 2 &&  & \vdash ( ( ( \mw{\varphi} \wedge \mw{\psi} \wedge \mw{\chi} )
      \leftrightarrow ( ( \mw{\varphi} \wedge \mw{\psi} )
      \wedge \mw{\chi} ) ) \rightarrow \mathbf{I} \lsem ( (
      \mw{\varphi} \wedge \mw{\psi} \wedge \mw{\chi} )
      \leftrightarrow ( ( \mw{\varphi} \wedge \mw{\psi} )
      \wedge \mw{\chi} ) ) \rsem
      )&(\mbox{\ref{eq:ax-bl.taut}})\notag
  \\ 3 &&  & \vdash \mathbf{I} \lsem ( ( \mw{\varphi} \wedge \mw{\psi} \wedge
      \mw{\chi} ) \leftrightarrow ( ( \mw{\varphi} \wedge
      \mw{\psi} ) \wedge \mw{\chi} ) )
      \rsem&(\mbox{\ref{eq:ax-mp}},1,2)\notag
  \\ 4 &&  & \vdash ( \mathbf{I} \lsem ( ( \mw{\varphi} \wedge \mw{\psi} \wedge
      \mw{\chi} ) \leftrightarrow ( ( \mw{\varphi} \wedge
      \mw{\psi} ) \wedge \mw{\chi} ) ) \rsem
      \rightarrow ( \mathbf{I} \lsem ( \mw{\varphi}
      \wedge \mw{\psi} \wedge \mw{\chi} ) \rsem
      \leftrightarrow \mathbf{I} \lsem ( ( \mw{\varphi}
      \wedge \mw{\psi} ) \wedge \mw{\chi} ) \rsem )
      )&(\mbox{\ref{eq:df-bl.bi}})\notag
  \\ 5 &&  & \vdash ( \mathbf{I} \lsem ( \mw{\varphi} \wedge \mw{\psi} \wedge
      \mw{\chi} ) \rsem \leftrightarrow \mathbf{I} \lsem (
      ( \mw{\varphi} \wedge \mw{\psi} ) \wedge \mw{\chi} )
      \rsem )&(\mbox{\ref{eq:ax-mp}},3,4)\notag
  \\ 6 &&  & \vdash ( \mathbf{I} \lsem ( ( \mw{\varphi} \wedge \mw{\psi} )
      \wedge \mw{\chi} ) \rsem \leftrightarrow ( \mathbf{I}
      \lsem ( \mw{\varphi} \wedge \mw{\psi} ) \rsem \wedge
      \mathbf{I} \lsem \mw{\chi} \rsem )
      )&(\mbox{\ref{eq:df-bl.an}})\notag
  \\ 7 &&  & \vdash ( \mathbf{I} \lsem ( \mw{\varphi} \wedge \mw{\psi} \wedge
      \mw{\chi} ) \rsem \leftrightarrow ( \mathbf{I} \lsem
      ( \mw{\varphi} \wedge \mw{\psi} ) \rsem \wedge
      \mathbf{I} \lsem \mw{\chi} \rsem )
      )&(\mbox{\ref{eq:bitri}},5,6)\notag
  \\ 8 &&  & \vdash ( \mathbf{I} \lsem ( \mw{\varphi} \wedge \mw{\psi} ) \rsem
      \leftrightarrow ( \mathbf{I} \lsem \mw{\varphi} \rsem
      \wedge \mathbf{I} \lsem \mw{\psi} \rsem )
      )&(\mbox{\ref{eq:df-bl.an}})\notag
  \\ 9 &&  & \vdash ( ( \mathbf{I} \lsem ( \mw{\varphi} \wedge \mw{\psi} )
      \rsem \wedge \mathbf{I} \lsem \mw{\chi} \rsem )
      \leftrightarrow ( ( \mathbf{I} \lsem \mw{\varphi}
      \rsem \wedge \mathbf{I} \lsem \mw{\psi} \rsem )
      \wedge \mathbf{I} \lsem \mw{\chi} \rsem )
      )&(\mbox{\ref{eq:anbi1i}},8)\notag
  \\ 10 &&  & \vdash ( \mathbf{I} \lsem ( \mw{\varphi} \wedge \mw{\psi} \wedge
      \mw{\chi} ) \rsem \leftrightarrow ( ( \mathbf{I}
      \lsem \mw{\varphi} \rsem \wedge \mathbf{I} \lsem
      \mw{\psi} \rsem ) \wedge \mathbf{I} \lsem \mw{\chi}
      \rsem ) )&(\mbox{\ref{eq:bitri}},7,9)\notag
  \\ 11 &&  & \vdash ( ( \mathbf{I} \lsem \mw{\varphi} \rsem \wedge \mathbf{I}
      \lsem \mw{\psi} \rsem \wedge \mathbf{I} \lsem
      \mw{\chi} \rsem ) \leftrightarrow ( ( \mathbf{I}
      \lsem \mw{\varphi} \rsem \wedge \mathbf{I} \lsem
      \mw{\psi} \rsem ) \wedge \mathbf{I} \lsem \mw{\chi}
      \rsem ) )&(\mbox{\ref{eq:df-3an}})\notag
  \\ 12 &&  & \vdash ( \mathbf{I} \lsem ( \mw{\varphi} \wedge \mw{\psi} \wedge
      \mw{\chi} ) \rsem \leftrightarrow ( \mathbf{I} \lsem
      \mw{\varphi} \rsem \wedge \mathbf{I} \lsem \mw{\psi}
      \rsem \wedge \mathbf{I} \lsem \mw{\chi} \rsem )
      )&(\mbox{\ref{eq:bitr4i}},10,11)\notag
\end{align}
\end{proof}

\begin{restatable}{metadefinition}{dfblor}~\blTitle{df-bl.or}~ Interpretation of disjunction.  Only
the sound direction is axiomatized: the converse
($\mathbf{I}\lsem\mw{\varphi}\vee\mw{\psi}\rsem \rightarrow
\mathbf{I}\lsem\mw{\varphi}\rsem\vee\mathbf{I}\lsem\mw{\psi}\rsem$) is false
in general (Lean4SFS/SFS.lean's dl\_or\_sound).
\begin{align}
\vdash ( ( \mathbf{I} \lsem \mw{\varphi} \rsem \vee \mathbf{I} \lsem \mw{\psi}
    \rsem ) \rightarrow \mathbf{I} \lsem ( \mw{\varphi} \vee \mw{\psi} ) \rsem
    )\label{eq:df-bl.or}\tag{df-bl.or}
\end{align}
\end{restatable}

\begin{restatable}{lemma}{bl3or}~\blTitle{bl.3or}~ Interpretation of disjunction, three wff.  Only the
sound direction is axiomatized: the converse is false in general, same as
df-bl.or itself (Lean4SFS/SFS.lean's dl\_3or\_sound).
\begin{align}
\vdash ( ( \mathbf{I} \lsem \mw{\varphi} \rsem \vee \mathbf{I} \lsem \mw{\psi}
    \rsem \vee \mathbf{I} \lsem \mw{\chi} \rsem ) \rightarrow \mathbf{I} \lsem
    ( \mw{\varphi} \vee \mw{\psi} \vee \mw{\chi} ) \rsem
    )\label{eq:bl.3or}\tag{bl.3or}
\end{align}
\end{restatable}

\begin{proof}
\begin{align}
  1 &&  & \vdash ( ( \mathbf{I} \lsem \mw{\varphi} \rsem \vee \mathbf{I}
      \lsem \mw{\psi} \rsem \vee \mathbf{I} \lsem \mw{\chi} \rsem )
      \leftrightarrow ( ( \mathbf{I} \lsem \mw{\varphi} \rsem \vee
      \mathbf{I} \lsem \mw{\psi} \rsem ) \vee \mathbf{I} \lsem \mw{\chi}
      \rsem ) )&(\mbox{\ref{eq:df-3or}})\notag
  \\ 2 &&  & \vdash ( ( \mathbf{I} \lsem \mw{\varphi} \rsem \vee
      \mathbf{I} \lsem \mw{\psi} \rsem \vee \mathbf{I} \lsem \mw{\chi}
      \rsem ) \rightarrow ( ( \mathbf{I} \lsem \mw{\varphi} \rsem \vee
      \mathbf{I} \lsem \mw{\psi} \rsem ) \vee \mathbf{I} \lsem
      \mw{\chi} \rsem ) )&(\mbox{\ref{eq:biimpi}},1)\notag
  \\ 3 &&  & \vdash ( ( \mathbf{I} \lsem \mw{\varphi} \rsem \vee
      \mathbf{I} \lsem \mw{\psi} \rsem ) \rightarrow \mathbf{I} \lsem (
      \mw{\varphi} \vee \mw{\psi} ) \rsem
      )&(\mbox{\ref{eq:df-bl.or}})\notag
  \\ 4 &&  & \vdash ( ( ( \mathbf{I} \lsem \mw{\varphi} \rsem \vee
      \mathbf{I} \lsem \mw{\psi} \rsem ) \vee \mathbf{I} \lsem
      \mw{\chi} \rsem ) \rightarrow ( \mathbf{I} \lsem ( \mw{\varphi}
      \vee \mw{\psi} ) \rsem \vee \mathbf{I} \lsem \mw{\chi} \rsem )
      )&(\mbox{\ref{eq:orim1i}},3)\notag
  \\ 5 &&  & \vdash ( ( \mathbf{I} \lsem \mw{\varphi} \rsem \vee
      \mathbf{I} \lsem \mw{\psi} \rsem \vee \mathbf{I} \lsem \mw{\chi}
      \rsem ) \rightarrow ( \mathbf{I} \lsem ( \mw{\varphi} \vee
      \mw{\psi} ) \rsem \vee \mathbf{I} \lsem \mw{\chi} \rsem )
      )&(\mbox{\ref{eq:syl}},2,4)\notag
  \\ 6 &&  & \vdash ( ( \mathbf{I} \lsem ( \mw{\varphi} \vee \mw{\psi}
      ) \rsem \vee \mathbf{I} \lsem \mw{\chi} \rsem ) \rightarrow
      \mathbf{I} \lsem ( ( \mw{\varphi} \vee \mw{\psi} ) \vee \mw{\chi}
      ) \rsem )&(\mbox{\ref{eq:df-bl.or}})\notag
  \\ 7 &&  & \vdash ( ( \mathbf{I} \lsem \mw{\varphi} \rsem \vee
      \mathbf{I} \lsem \mw{\psi} \rsem \vee \mathbf{I} \lsem \mw{\chi}
      \rsem ) \rightarrow \mathbf{I} \lsem ( ( \mw{\varphi} \vee
      \mw{\psi} ) \vee \mw{\chi} ) \rsem
      )&(\mbox{\ref{eq:syl}},5,6)\notag
  \\ 8 &&  & \vdash ( ( \mw{\varphi} \vee \mw{\psi} \vee \mw{\chi} )
      \leftrightarrow ( ( \mw{\varphi} \vee \mw{\psi} ) \vee \mw{\chi}
      ) )&(\mbox{\ref{eq:df-3or}})\notag
  \\ 9 &&  & \vdash ( ( ( \mw{\varphi} \vee \mw{\psi} ) \vee \mw{\chi}
      ) \rightarrow ( \mw{\varphi} \vee \mw{\psi} \vee \mw{\chi} )
      )&(\mbox{\ref{eq:biimpri}},8)\notag
  \\ 10 &&  & \vdash ( ( ( ( \mw{\varphi} \vee \mw{\psi} ) \vee
      \mw{\chi} ) \rightarrow ( \mw{\varphi} \vee \mw{\psi} \vee
      \mw{\chi} ) ) \rightarrow ( \mathbf{I} \lsem ( ( \mw{\varphi}
      \vee \mw{\psi} ) \vee \mw{\chi} ) \rsem \rightarrow \mathbf{I}
      \lsem ( \mw{\varphi} \vee \mw{\psi} \vee \mw{\chi} ) \rsem )
      )&(\mbox{\ref{eq:bl.mono}})\notag
  \\ 11 &&  & \vdash ( \mathbf{I} \lsem ( ( \mw{\varphi} \vee
      \mw{\psi} ) \vee \mw{\chi} ) \rsem \rightarrow \mathbf{I} \lsem
      ( \mw{\varphi} \vee \mw{\psi} \vee \mw{\chi} ) \rsem
      )&(\mbox{\ref{eq:ax-mp}},9,10)\notag
  \\ 12 &&  & \vdash ( ( \mathbf{I} \lsem \mw{\varphi} \rsem \vee
      \mathbf{I} \lsem \mw{\psi} \rsem \vee \mathbf{I} \lsem \mw{\chi}
      \rsem ) \rightarrow \mathbf{I} \lsem ( \mw{\varphi} \vee
      \mw{\psi} \vee \mw{\chi} ) \rsem
      )&(\mbox{\ref{eq:syl}},7,11)\notag
\end{align}
\end{proof}

\begin{restatable}{metadefinition}{dfblnot}~\blTitle{df-bl.not}~ Interpretation of complement.  Only
the sound direction is axiomatized: the converse
($\lnot\mathbf{I}\lsem\mw{\varphi}\rsem \rightarrow
\mathbf{I}\lsem\lnot\mw{\varphi}\rsem$) is false in general (Lean4SFS/
SFS.lean's dl\_not\_sound).
\begin{align}
\vdash ( \mathbf{I} \lsem \lnot \mw{\varphi} \rsem \rightarrow \lnot
    \mathbf{I} \lsem \mw{\varphi} \rsem )\label{eq:df-bl.not}\tag{df-bl.not}
\end{align}
\end{restatable}

\begin{restatable}{metadefinition}{dfblal}~\blTitle{df-bl.al}~ Exportation of universal quantification
from interpretation $\mathbf{I}$.
\begin{align}
\vdash ( \mathbf{I} \lsem \forall \msc{x} \in \mc{A}~\mw{\varphi} \rsem
    \leftrightarrow \forall \msc{x} \in \mc{A} \mathbf{I} \lsem \mw{\varphi}
    \rsem )\label{eq:df-bl.al}\tag{df-bl.al}
\end{align}
\end{restatable}

\begin{restatable}{lemma}{blmonobi}~\blTitle{bl.monobi}~ $\mathbf{I}$~respects logical equivalence -- the
iff-congruence counterpart of bl.mono, same derivation pattern
(ax-bl.taut necessitation + df-bl.bi's sound direction).
\begin{align}
\vdash ( ( \mw{\varphi} \leftrightarrow \mw{\psi} ) \rightarrow ( \mathbf{I}
    \lsem \mw{\varphi} \rsem \leftrightarrow \mathbf{I} \lsem \mw{\psi} \rsem
    ) )\label{eq:bl.monobi}\tag{bl.monobi}
\end{align}
\end{restatable}

\begin{proof}
\begin{align}
  1 &&  & \vdash ( ( \mw{\varphi} \leftrightarrow \mw{\psi} ) \rightarrow
      \mathbf{I} \lsem ( \mw{\varphi} \leftrightarrow \mw{\psi} ) \rsem
      )&(\mbox{\ref{eq:ax-bl.taut}})\notag
  \\ 2 &&  & \vdash ( \mathbf{I} \lsem ( \mw{\varphi} \leftrightarrow
      \mw{\psi} ) \rsem \rightarrow ( \mathbf{I} \lsem \mw{\varphi}
      \rsem \leftrightarrow \mathbf{I} \lsem \mw{\psi} \rsem )
      )&(\mbox{\ref{eq:df-bl.bi}})\notag
  \\ 3 &&  & \vdash ( ( \mw{\varphi} \leftrightarrow \mw{\psi} )
      \rightarrow ( \mathbf{I} \lsem \mw{\varphi} \rsem \leftrightarrow
      \mathbf{I} \lsem \mw{\psi} \rsem )
      )&(\mbox{\ref{eq:syl}},1,2)\notag
\end{align}
\end{proof}

\begin{restatable}{lemma}{blal}~\blTitle{bl.al}~ Unrestricted version of df-bl.al, via ralv
($\forall \msc{x} \in \,\mathrm{V}~\mw{\varphi} \leftrightarrow \forall
\msc{x}~\mw{\varphi}$).
\begin{align}
\vdash ( \mathbf{I} \lsem \forall \msc{x} \mw{\varphi} \rsem \leftrightarrow
    \forall \msc{x} \mathbf{I} \lsem \mw{\varphi} \rsem
    )\label{eq:bl.al}\tag{bl.al}
\end{align}
\end{restatable}

\begin{proof}
\begin{align}
  1 &&  & \vdash ( \forall \msc{x} \in \,\mathrm{V} \mw{\varphi}
      \leftrightarrow \forall \msc{x} \mw{\varphi}
      )&(\mbox{\ref{eq:ralv}})\notag
  \\ 2 &&  & \vdash ( \forall \msc{x} \mw{\varphi} \leftrightarrow
      \forall \msc{x} \in \,\mathrm{V} \mw{\varphi}
      )&(\mbox{\ref{eq:bicomi}},1)\notag
  \\ 3 &&  & \vdash ( ( \forall \msc{x} \mw{\varphi} \leftrightarrow
      \forall \msc{x} \in \,\mathrm{V} \mw{\varphi} ) \rightarrow (
      \mathbf{I} \lsem \forall \msc{x} \mw{\varphi} \rsem
      \leftrightarrow \mathbf{I} \lsem \forall \msc{x} \in \,\mathrm{V}
      \mw{\varphi} \rsem ) )&(\mbox{\ref{eq:bl.monobi}})\notag
  \\ 4 &&  & \vdash ( \mathbf{I} \lsem \forall \msc{x} \mw{\varphi}
      \rsem \leftrightarrow \mathbf{I} \lsem \forall \msc{x} \in
      \,\mathrm{V} \mw{\varphi} \rsem
      )&(\mbox{\ref{eq:ax-mp}},2,3)\notag
  \\ 5 &&  & \vdash ( \mathbf{I} \lsem \forall \msc{x} \in \,\mathrm{V}
      \mw{\varphi} \rsem \leftrightarrow \forall \msc{x} \in
      \,\mathrm{V} \mathbf{I} \lsem \mw{\varphi} \rsem
      )&(\mbox{\ref{eq:df-bl.al}})\notag
  \\ 6 &&  & \vdash ( \forall \msc{x} \in \,\mathrm{V} \mathbf{I} \lsem
      \mw{\varphi} \rsem \leftrightarrow \forall \msc{x} \mathbf{I}
      \lsem \mw{\varphi} \rsem )&(\mbox{\ref{eq:ralv}})\notag
  \\ 7 &&  & \vdash ( \mathbf{I} \lsem \forall \msc{x} \in \,\mathrm{V}
      \mw{\varphi} \rsem \leftrightarrow \forall \msc{x} \mathbf{I}
      \lsem \mw{\varphi} \rsem )&(\mbox{\ref{eq:bitri}},5,6)\notag
  \\ 8 &&  & \vdash ( \mathbf{I} \lsem \forall \msc{x} \mw{\varphi}
      \rsem \leftrightarrow \forall \msc{x} \mathbf{I} \lsem
      \mw{\varphi} \rsem )&(\mbox{\ref{eq:bitri}},4,7)\notag
\end{align}
\end{proof}

\begin{restatable}{lemma}{blnf}~\blTitle{bl.nf}~ $\mathbf{I}\lsem\mw{\varphi}\rsem$ is vacuously quantifiable
whenever $\mw{\varphi}$ is: needed so bl.ex's witness-elimination step
(rexlimi) has the $\Finv$ hypothesis it needs on the boldI-wrapped side,
not just the raw side (nfre1).
\begin{align}
\vdash \Finv \msc{x} \mw{\varphi}\quad\Rightarrow\quad \vdash \Finv \msc{x}
    \mathbf{I} \lsem \mw{\varphi} \rsem\label{eq:bl.nf}\tag{bl.nf}
\end{align}
\end{restatable}

\begin{proof}
\begin{align}
  1 &&  & \vdash ( \mathbf{I} \lsem \forall \msc{x} \mw{\varphi} \rsem
      \leftrightarrow \forall \msc{x} \mathbf{I} \lsem \mw{\varphi} \rsem
      )&(\mbox{\ref{eq:bl.al}})\notag
  \\ 2 &&  & \vdash \Finv \msc{x}
      \mw{\varphi}&\text{Hyp~nf.1}\notag
  \\ 3 &&  & \vdash ( \forall \msc{x} \mw{\varphi} \leftrightarrow
      \mw{\varphi} )&(\mbox{\ref{eq:19.3}},2)\notag
  \\ 4 &&  & \vdash ( ( \forall \msc{x} \mw{\varphi} \leftrightarrow
      \mw{\varphi} ) \rightarrow ( \mathbf{I} \lsem \forall \msc{x}
      \mw{\varphi} \rsem \leftrightarrow \mathbf{I} \lsem \mw{\varphi}
      \rsem ) )&(\mbox{\ref{eq:bl.monobi}})\notag
  \\ 5 &&  & \vdash ( \mathbf{I} \lsem \forall \msc{x} \mw{\varphi}
      \rsem \leftrightarrow \mathbf{I} \lsem \mw{\varphi} \rsem
      )&(\mbox{\ref{eq:ax-mp}},3,4)\notag
  \\ 6 &&  & \vdash ( \forall \msc{x} \mathbf{I} \lsem \mw{\varphi}
      \rsem \leftrightarrow \mathbf{I} \lsem \mw{\varphi} \rsem
      )&(\mbox{\ref{eq:bitr3i}},1,5)\notag
  \\ 7 &&  & \vdash ( \mathbf{I} \lsem \mw{\varphi} \rsem \rightarrow
      \forall \msc{x} \mathbf{I} \lsem \mw{\varphi} \rsem
      )&(\mbox{\ref{eq:biimpri}},6)\notag
  \\ 8 &&  & \vdash \forall \msc{x} ( \mathbf{I} \lsem \mw{\varphi}
      \rsem \rightarrow \forall \msc{x} \mathbf{I} \lsem \mw{\varphi}
      \rsem )&(\mbox{\ref{eq:ax-gen}},7)\notag
  \\ 9 &&  & \vdash ( \forall \msc{x} ( \mathbf{I} \lsem \mw{\varphi}
      \rsem \rightarrow \forall \msc{x} \mathbf{I} \lsem \mw{\varphi}
      \rsem ) \rightarrow \Finv \msc{x} \mathbf{I} \lsem \mw{\varphi}
      \rsem )&(\mbox{\ref{eq:nf5-1}})\notag
  \\ 10 &&  & \vdash \Finv \msc{x} \mathbf{I} \lsem \mw{\varphi}
      \rsem&(\mbox{\ref{eq:ax-mp}},8,9)\notag
\end{align}
\end{proof}

\begin{restatable}{lemma}{bldfrex2}~\blTitle{bl.dfrex2}~ Define existential quantification in
interpretation $\mathbf{I}$.
\begin{align}
\vdash ( \mathbf{I} \lsem \exists \msc{x} \in \mc{A}~\mw{\varphi} \rsem
    \leftrightarrow \mathbf{I} \lsem \lnot \forall \msc{x} \in \mc{A} \lnot
    \mw{\varphi} \rsem )\label{eq:bl.dfrex2}\tag{bl.dfrex2}
\end{align}
\end{restatable}

\begin{proof}
\begin{align}
  1 &&  & \vdash ( \exists \msc{x} \in \mc{A}~\mw{\varphi} \leftrightarrow
      \lnot \forall \msc{x} \in \mc{A} \lnot \mw{\varphi}
      )&(\mbox{\ref{eq:dfrex2}})\notag
  \\ 2 &&  & \vdash ( ( \exists \msc{x} \in \mc{A}~\mw{\varphi} \leftrightarrow
      \lnot \forall \msc{x} \in \mc{A} \lnot \mw{\varphi} )
      \rightarrow \mathbf{I} \lsem ( \exists \msc{x} \in
      \mc{A}~\mw{\varphi} \leftrightarrow \lnot \forall
      \msc{x} \in \mc{A} \lnot \mw{\varphi} ) \rsem
      )&(\mbox{\ref{eq:ax-bl.taut}})\notag
  \\ 3 &&  & \vdash \mathbf{I} \lsem ( \exists \msc{x} \in \mc{A}~\mw{\varphi}
      \leftrightarrow \lnot \forall \msc{x} \in \mc{A}
      \lnot \mw{\varphi} )
      \rsem&(\mbox{\ref{eq:ax-mp}},1,2)\notag
  \\ 4 &&  & \vdash ( \mathbf{I} \lsem ( \exists \msc{x} \in \mc{A}
      \mw{\varphi} \leftrightarrow \lnot \forall \msc{x}
      \in \mc{A} \lnot \mw{\varphi} ) \rsem \rightarrow
      ( \mathbf{I} \lsem \exists \msc{x} \in \mc{A}
      \mw{\varphi} \rsem \leftrightarrow \mathbf{I} \lsem
      \lnot \forall \msc{x} \in \mc{A} \lnot \mw{\varphi}
      \rsem ) )&(\mbox{\ref{eq:df-bl.bi}})\notag
  \\ 5 &&  & \vdash ( \mathbf{I} \lsem \exists \msc{x} \in \mc{A}~\mw{\varphi}
      \rsem \leftrightarrow \mathbf{I} \lsem \lnot \forall
      \msc{x} \in \mc{A} \lnot \mw{\varphi} \rsem
      )&(\mbox{\ref{eq:ax-mp}},3,4)\notag
\end{align}
\end{proof}

\begin{restatable}{lemma}{blexraw}~\blTitle{bl.exraw}~ Raw (non-modal) restricted-existential-introduction
fact, unconditional: the ordinary logic behind bl.ex, before any $\mathbf{I}$ is
involved.
\begin{align}
\vdash ( ( \msc{x} \in \mc{A} \wedge \mw{\varphi} ) \rightarrow \exists
    \msc{x} \in \mc{A}~\mw{\varphi} )\label{eq:bl.exraw}\tag{bl.exraw}
\end{align}
\end{restatable}

\begin{proof}
\begin{align}
  1 &&  & \vdash ( ( \msc{x} \in \mc{A} \wedge \mw{\varphi} ) \rightarrow
      \exists \msc{x} ( \msc{x} \in \mc{A} \wedge \mw{\varphi} )
      )&(\mbox{\ref{eq:19.8a}})\notag
  \\ 2 &&  & \vdash ( \exists \msc{x} \in \mc{A}~\mw{\varphi}
      \leftrightarrow \exists \msc{x} ( \msc{x} \in \mc{A} \wedge
      \mw{\varphi} ) )&(\mbox{\ref{eq:df-rex}})\notag
  \\ 3 &&  & \vdash ( \exists \msc{x} ( \msc{x} \in \mc{A} \wedge
      \mw{\varphi} ) \rightarrow \exists \msc{x} \in \mc{A}
      \mw{\varphi} )&(\mbox{\ref{eq:biimpri}},2)\notag
  \\ 4 &&  & \vdash ( ( \msc{x} \in \mc{A} \wedge \mw{\varphi} )
      \rightarrow \exists \msc{x} \in \mc{A}~\mw{\varphi}
      )&(\mbox{\ref{eq:syl}},1,3)\notag
\end{align}
\end{proof}

\begin{restatable}{lemma}{blexlemzero}~\blTitle{bl.exlem0}~ Helper for bl.ex.  Built from bl.exraw,
bl.mono, bl.ty (types are constant under interpretation -- one of this
file's own pre-existing unproven \$p placeholders, so this lemma inherits
that gap rather than introducing a new one), and df-bl.an (sound both ways).
\begin{align}
\vdash ( ( \msc{x} \in \mc{A} \wedge \mathbf{I} \lsem \mw{\varphi} \rsem )
    \rightarrow \mathbf{I} \lsem \exists \msc{x} \in \mc{A}~\mw{\varphi} \rsem
    )\label{eq:bl.exlem0}\tag{bl.exlem0}
\end{align}
\end{restatable}

\begin{proof}
\begin{align}
  1 &&  & \vdash ( \msc{x} \in \mc{A} \rightarrow \mathbf{I} \lsem \msc{x}
      \in \mc{A} \rsem )&(\mbox{\ref{eq:bl.ty}})\notag
  \\ 2 &&  & \vdash ( ( \msc{x} \in \mc{A} \wedge \mathbf{I} \lsem
      \mw{\varphi} \rsem ) \rightarrow ( \mathbf{I} \lsem \msc{x} \in
      \mc{A} \rsem \wedge \mathbf{I} \lsem \mw{\varphi} \rsem )
      )&(\mbox{\ref{eq:anim1i}},1)\notag
  \\ 3 &&  & \vdash ( \mathbf{I} \lsem ( \msc{x} \in \mc{A} \wedge
      \mw{\varphi} ) \rsem \leftrightarrow ( \mathbf{I} \lsem \msc{x}
      \in \mc{A} \rsem \wedge \mathbf{I} \lsem \mw{\varphi} \rsem )
      )&(\mbox{\ref{eq:df-bl.an}})\notag
  \\ 4 &&  & \vdash ( ( \mathbf{I} \lsem \msc{x} \in \mc{A} \rsem
      \wedge \mathbf{I} \lsem \mw{\varphi} \rsem ) \rightarrow
      \mathbf{I} \lsem ( \msc{x} \in \mc{A} \wedge \mw{\varphi} ) \rsem
      )&(\mbox{\ref{eq:biimpri}},3)\notag
  \\ 5 &&  & \vdash ( ( \msc{x} \in \mc{A} \wedge \mathbf{I} \lsem
      \mw{\varphi} \rsem ) \rightarrow \mathbf{I} \lsem ( \msc{x} \in
      \mc{A} \wedge \mw{\varphi} ) \rsem
      )&(\mbox{\ref{eq:syl}},2,4)\notag
  \\ 6 &&  & \vdash ( ( \msc{x} \in \mc{A} \wedge \mw{\varphi} )
      \rightarrow \exists \msc{x} \in \mc{A}~\mw{\varphi}
      )&(\mbox{\ref{eq:bl.exraw}})\notag
  \\ 7 &&  & \vdash ( ( ( \msc{x} \in \mc{A} \wedge \mw{\varphi} )
      \rightarrow \exists \msc{x} \in \mc{A}~\mw{\varphi} ) \rightarrow
      ( \mathbf{I} \lsem ( \msc{x} \in \mc{A} \wedge \mw{\varphi} )
      \rsem \rightarrow \mathbf{I} \lsem \exists \msc{x} \in \mc{A}
      \mw{\varphi} \rsem ) )&(\mbox{\ref{eq:bl.mono}})\notag
  \\ 8 &&  & \vdash ( \mathbf{I} \lsem ( \msc{x} \in \mc{A} \wedge
      \mw{\varphi} ) \rsem \rightarrow \mathbf{I} \lsem \exists \msc{x}
      \in \mc{A}~\mw{\varphi} \rsem )&(\mbox{\ref{eq:ax-mp}},6,7)\notag
  \\ 9 &&  & \vdash ( ( \msc{x} \in \mc{A} \wedge \mathbf{I} \lsem
      \mw{\varphi} \rsem ) \rightarrow \mathbf{I} \lsem \exists \msc{x}
      \in \mc{A}~\mw{\varphi} \rsem )&(\mbox{\ref{eq:syl}},5,8)\notag
\end{align}
\end{proof}

\begin{restatable}{lemma}{blexlem}~\blTitle{bl.exlem}~ Exported form of bl.exlem0.
\begin{align}
\vdash ( \msc{x} \in \mc{A} \rightarrow ( \mathbf{I} \lsem \mw{\varphi} \rsem
    \rightarrow \mathbf{I} \lsem \exists \msc{x} \in \mc{A}~\mw{\varphi} \rsem
    ) )\label{eq:bl.exlem}\tag{bl.exlem}
\end{align}
\end{restatable}

\begin{proof}
\begin{align}
  1 &&  & \vdash ( ( \msc{x} \in \mc{A} \wedge \mathbf{I} \lsem
      \mw{\varphi} \rsem ) \rightarrow \mathbf{I} \lsem \exists \msc{x}
      \in \mc{A}~\mw{\varphi} \rsem )&(\mbox{\ref{eq:bl.exlem0}})\notag
  \\ 2 &&  & \vdash ( \msc{x} \in \mc{A} \rightarrow ( \mathbf{I} \lsem
      \mw{\varphi} \rsem \rightarrow \mathbf{I} \lsem \exists \msc{x}
      \in \mc{A}~\mw{\varphi} \rsem ) )&(\mbox{\ref{eq:ex}},1)\notag
\end{align}
\end{proof}

\begin{restatable}{lemma}{blex}~\blTitle{bl.ex}~ Exportation of existential quantification
from interpretation $\mathbf{I}$.  Only the sound direction is axiomatized:
the converse (the classic invalid $\exists$-to-$\forall$-swap-under-$\lnot$
move -- different instants may need different witnesses) is false in
general (Lean4SFS/SFS.lean's dl\_ex\_sound).
\begin{align}
\vdash ( \exists \msc{x} \in \mc{A} \mathbf{I} \lsem \mw{\varphi} \rsem
    \rightarrow \mathbf{I} \lsem \exists \msc{x} \in \mc{A}~\mw{\varphi} \rsem
    )\label{eq:bl.ex}\tag{bl.ex}
\end{align}
\end{restatable}

\begin{proof}
\begin{align}
  1 &&  & \vdash \Finv \msc{x} \exists \msc{x} \in \mc{A}
      \mw{\varphi}&(\mbox{\ref{eq:nfre1}})\notag
  \\ 2 &&  & \vdash \Finv \msc{x} \mathbf{I} \lsem \exists \msc{x} \in
      \mc{A} \mw{\varphi} \rsem&(\mbox{\ref{eq:bl.nf}},1)\notag
  \\ 3 &&  & \vdash ( \msc{x} \in \mc{A} \rightarrow ( \mathbf{I} \lsem
      \mw{\varphi} \rsem \rightarrow \mathbf{I} \lsem \exists \msc{x}
      \in \mc{A}~\mw{\varphi} \rsem )
      )&(\mbox{\ref{eq:bl.exlem}})\notag
  \\ 4 &&  & \vdash ( \exists \msc{x} \in \mc{A} \mathbf{I} \lsem
      \mw{\varphi} \rsem \rightarrow \mathbf{I} \lsem \exists \msc{x}
      \in \mc{A}~\mw{\varphi} \rsem
      )&(\mbox{\ref{eq:rexlimi}},2,3)\notag
\end{align}
\end{proof}

\begin{restatable}{metadefinition}{dfbladd}~\blTitle{df-bl.add}~ Interpretation of addition.
\begin{align}
\vdash \mathbf{I} \lsem ( \mc{A} + \mc{B} ) \rsem = ( \mathbf{I} \lsem \mc{A}
    \rsem + \mathbf{I} \lsem \mc{B} \rsem )\label{eq:df-bl.add}\tag{df-bl.add}
\end{align}
\end{restatable}

\begin{restatable}{metadefinition}{dfblsub}~\blTitle{df-bl.sub}~ Interpretation of subtraction.
\begin{align}
\vdash \mathbf{I} \lsem ( \mc{A} \minus \mc{B} ) \rsem = ( \mathbf{I} \lsem
    \mc{A} \rsem \minus \mathbf{I} \lsem \mc{B} \rsem
    )\label{eq:df-bl.sub}\tag{df-bl.sub}
\end{align}
\end{restatable}

\begin{restatable}{metadefinition}{dfblmul}~\blTitle{df-bl.mul}~ Interpretation of multiplication.
\begin{align}
\vdash \mathbf{I} \lsem ( \mc{A} \times \mc{B} ) \rsem = ( \mathbf{I} \lsem
    \mc{A} \rsem \times \mathbf{I} \lsem \mc{B} \rsem
    )\label{eq:df-bl.mul}\tag{df-bl.mul}
\end{align}
\end{restatable}

\begin{restatable}{metadefinition}{dfbldiv}~\blTitle{df-bl.div}~ Interpretation of division.
\begin{align}
\vdash \mathbf{I} \lsem ( \mc{A} / \mc{B} ) \rsem = ( \mathbf{I} \lsem \mc{A}
    \rsem / \mathbf{I} \lsem \mc{B} \rsem )\label{eq:df-bl.div}\tag{df-bl.div}
\end{align}
\end{restatable}

\begin{restatable}{metadefinition}{dfblum}~\blTitle{df-bl.um}~ Interpretation of negation.
\begin{align}
\vdash \mathbf{I} \lsem \textrm{-} \mc{A} \rsem = \textrm{-} \mathbf{I} \lsem
    \mc{A} \rsem\label{eq:df-bl.um}\tag{df-bl.um}
\end{align}
\end{restatable}

\begin{restatable}{metadefinition}{dfbleq}~\blTitle{df-bl.eq}~ Interpretation of equality.
\begin{align}
\vdash ( \mathbf{I} \lsem \mc{A} = \mc{B} \rsem \leftrightarrow \mathbf{I}
    \lsem \mc{A} \rsem = \mathbf{I} \lsem \mc{B} \rsem
    )\label{eq:df-bl.eq}\tag{df-bl.eq}
\end{align}
\end{restatable}

\begin{restatable}{metadefinition}{dfbllt}~\blTitle{df-bl.lt}~ Interpretation of less-than.
\begin{align}
\vdash ( \mathbf{I} \lsem \mc{A} < \mc{B} \rsem \leftrightarrow \mathbf{I}
    \lsem \mc{A} \rsem < \mathbf{I} \lsem \mc{B} \rsem
    )\label{eq:df-bl.lt}\tag{df-bl.lt}
\end{align}
\end{restatable}

\begin{restatable}{metadefinition}{dfblam}~\blTitle{df-bl.am}~ Interpretation of at-most.
\begin{align}
\vdash ( \mathbf{I} \lsem \mc{A} \le \mc{B} \rsem \leftrightarrow \mathbf{I}
    \lsem \mc{A} \rsem \le \mathbf{I} \lsem \mc{B} \rsem
    )\label{eq:df-bl.am}\tag{df-bl.am}
\end{align}
\end{restatable}

